\section{Polynomial boundary bounds without separating the zeros}
\label{sec:trace-frequency-boundary}

The local proof AC-T39 can be quantified without dividing by
$L'(\rho,\chi)$ or by a distance between distinct zeros.
The frequency cutoff below is allowed to depend on the height.
Its support and its plateau have fixed sizes.

\begin{theorem}[AC-T40: a uniformly controlled local inverse]
\label{AC-T40}
Fix a nonprincipal primitive Dirichlet character $\chi$ and
assume A.  Put $H=2+|T|$, with $|T|\geq10$.
For any fixed nonnegative integers $J,k$ there is a real
$a_T\in C_c^\infty((T-1,T+1))$ such that
\[
 0\leq a_T\leq1,\qquad a_T=1\ \text{on }[T-\tfrac12,T+\tfrac12].
                                                               \tag{FB.1}
\]
There is a list of $N\leq C_\chi\log H$ zeros
$\rho_j=\tfrac12+i\gamma_j$, counted with multiplicity,
containing every zero at an ordinate in the support of $a_T$.
Give every pole the prescription
$(\tau-\gamma_j-i0)^{-1}$, and let
\[
 b_T(\tau)=\frac{a_T(\tau)}{L(\tfrac12+i\tau,\chi)},\qquad
 \kappa_T(t)=\frac1{2\pi}\int_{\mathbb R}b_T(\tau)e^{-i\tau t}\,d\tau
                                                               \tag{FB.2}
\]
in this distributional sense.
For each fixed $A_0\geq0$ there are constants $C,c$, depending
only on $\chi,J,k,A_0$, such that
\[
 \sum_{\ell=0}^k\int_{-A_0}^{\infty}(1+|t|)^J
       \bigl(|\kappa_T^{(\ell)}(t)|
                    +|\kappa_T^{(\ell)}(t)|^2\bigr)\,dt
                   \leq C H^c.                              \tag{FB.3}
\]
The constants do not depend on multiplicities or zero gaps.
The same statement holds with an additional analytic numerator
$w_T(\tau)$, provided that on a fixed complex neighborhood of
$[T-1,T+1]$ its modulus is bounded by $C_0H^{c_0}$;
the constants may then depend on that neighborhood and on
$C_0,c_0$.

Let $S_L=1_{[-L,L]}$, $L\geq1$, and let $W_L$ be the
orthogonal projection in $S_LL^2(\mathbb R)$ onto the
orthogonal complement of all the zero jets
$S_Le^{-i\gamma t}t^j$ belonging to the listed zeros.
For this same cutoff, for every fixed $\ell\leq k$,
\[
 \|C_{b_T}W_L\|_{\mathcal S_2(L^2,H^\ell)}
                 \leq C H^c(1+L)^{N+2}.                     \tag{FB.4}
\]
Here $C_{b_T}$ uses the positive-sign Fourier transform,
as in AC-T39.  In (FB.4) the exponent of $1+L$ may be
increased by a fixed constant without affecting the claim.
In particular its logarithm is bounded by
$C\log H+(C\log H+2)\log(1+L)$.
\end{theorem}
\begin{proof}
\proofstep{1. A zero-free factor with polynomial bounds.}
Set $z_0=2+iT$.  On $|s-z_0|\leq4$,
\[
          |L(s,\chi)|\leq C_\chi H^{c_\chi},
                   \qquad |L(z_0,\chi)|\geq c_\chi'>0.
                                                               \tag{FB.5}
\]
For the lower bound, absolute Euler products on $\Re s=2$
give $|L(2+iT,\chi)|^{-1}\leq\prod_p(1+p^{-2})$.
For the upper bound one may use the exact finite decomposition
$L(s,\chi)=q^{-s}\sum_{a=1}^q\chi(a)\zeta(s,a/q)$.
The Euler--Maclaurin representation DLMF (25.11.7),
with $n=2$, has a bounded periodic Bernoulli remainder
integrated against $(x+a/q)^{-s-5}$.
It is absolutely bounded uniformly on $-2\leq\Re s\leq6$;
its prefactor is a polynomial of degree five in $s$.
The other terms have the same polynomial bound, since
$|\Im s|\geq6$ here.  This proves (FB.5).
The primary formula and its integration-by-parts derivation
are recorded at
\url{https://dlmf.nist.gov/25.11.E7};
this is a use of that known formula, not a new continuation
theorem.

Choose $r\in(11/4,3)$ with no zero on $|s-z_0|=r$.
Jensen's formula, on any zero-free boundary circle of radius
between $7/2$ and $4$, and (FB.5) give
\[
 \#\{\rho:|\rho-z_0|<r\}\log(7/6)\leq C_\chi\log H.
                                                               \tag{FB.6}
\]
Indeed each of those zeros contributes at least
$\log((7/2)/3)=\log(7/6)$ to the Jensen sum, whereas
the boundary mean of $\log|L|-\log|L(z_0)|$ is
$O_\chi(\log H)$.
List these $N$ zeros with multiplicity.
There are no trivial zeros in this disk, because $|T|\geq10$.
Theorem A places them on $\Re s=1/2$.
In particular $|\rho_j-z_0|\geq3/2$.

Write $z=s-z_0$, $\alpha_j=\rho_j-z_0$, and put
\[
 B_T(s)=\prod_{j=1}^N
       \frac{r(z-\alpha_j)}{r^2-\overline{\alpha_j}z},
                    \qquad H_T(s)=\frac{L(s,\chi)}{B_T(s)}.
                                                               \tag{FB.7}
\]
The denominators of the factors have their zeros outside
the disk.  Each factor has modulus one on its boundary.
Thus $H_T$ is analytic and zero-free on the closed disk,
and $|H_T|\leq C_\chi H^{c_\chi}$ there by the maximum
principle.  At its center,
$|H_T(z_0)|\geq|L(z_0,\chi)|\geq c_\chi'$.
Choose $M_0=C_\chi''\log H$ above the upper bound for
$\log|H_T|$.  The positive harmonic function
$M_0-\log|H_T|$ has value $O_\chi(\log H)$ at the center.
The Poisson kernel, or its immediate bound
$(r+\rho)/(r-\rho)$ on a concentric radius $\rho$,
therefore shows on $|s-z_0|\leq9/4$ that
\[
                   |H_T(s)|+|H_T(s)|^{-1}\leq C H^c.
                                                               \tag{FB.8}
\]
All radii are fixed, so these constants are uniform in the
choice of $r$ and in all zero positions.
Every $s=1/2+i\tau$ with $|\tau-T|\leq1$ has distance
at most $\sqrt{13}/2<9/4$ from $z_0$.
Cauchy's formula on a fixed smaller neighborhood gives
\[
 \left|\frac{d^m}{d\tau^m}
                H_T(\tfrac12+i\tau)^{-1}\right|
                     \leq C H^c C_1^m m!\quad(m\geq0).
                                                               \tag{FB.9}
\]
The optional analytic numerator satisfies the same
derivative bound with its stated additional power of $H$.

\proofstep{2. A cutoff with a controlled number of derivatives.}
Fix a nonnegative $C_c^\infty$ probability density $\eta$
supported in $(-1/16,1/16)$.
For the integer
\[
                         M=2N+2J+2k+12
\]
let $h_M=8M\,1_{[-1/(16M),1/(16M)]}$, and define
\[
 a_T(T+u)=
       \bigl(1_{[-3/4,3/4]}*h_M^{*M}*\eta\bigr)(u).
                                                               \tag{FB.10}
\]
The total smoothing support has radius at most $1/8$.
Hence this function is one on $[-5/8,5/8]$, zero outside
$[-7/8,7/8]$, and takes values in $[0,1]$, proving (FB.1).
For $0\leq m\leq M$, differentiate $m$ distinct box
factors.  Each derivative is a signed measure of total
variation $16M$, while every remaining smoothing factor
has mass one.  Young's inequality gives
\[
       \|a_T^{(m)}\|_\infty+\|a_T^{(m)}\|_2
                             \leq C(16M)^m.                  \tag{FB.11}
\]
Smoothness follows from the final convolution with $\eta$.

Put $f_T(\tau)=a_T(\tau)/H_T(1/2+i\tau)$, and let
$k_T$ be its inverse Fourier transform.
Leibniz's rule and (FB.9)--(FB.11), using
$m!\leq M^m$ for $m\leq M$, imply
\[
             \|f_T^{(m)}\|_2\leq C H^c(C_2M)^m
                          \quad(0\leq m\leq M).              \tag{FB.12}
\]
For any fixed $\ell\leq k+1$ the same bound, with an
additional fixed polynomial in $H,M$, holds for
derivatives of $( -i\tau)^\ell f_T(\tau)$.
The support has $|\tau|\leq|T|+1$; the finitely many
derivatives of $\tau^\ell$ prove this assertion directly.
Weighted Cauchy--Schwarz followed by Plancherel now gives,
for $0\leq m\leq M-1$,
\[
 \int_{\mathbb R}(1+|t|)^m|k_T^{(\ell)}(t)|\,dt
                  \leq C_\ell H^{c_\ell}
                                     (C_\ell M)^{m+C_\ell}.
                                                               \tag{FB.13}
\]
For example apply Cauchy--Schwarz with weight
$(1+t^2)^{-1}$ and bound the two $L^2$ norms of
$t^m k_T^{(\ell)}$ and $t^{m+1}k_T^{(\ell)}$ by
(FB.12).  Expanding $(1+|t|)^m$ only changes
$C_\ell^m$.  Fixed powers of $M$ will always be absorbed
in a power of $H$, since $M=O_{\chi,J,k}(\log H)$.
This argument also proves (FB.12)--(FB.13) with $w_T$.

\proofstep{3. The inverse zero product has no gap denominators.}
An individual inverse factor in (FB.7) is
\[
 \frac1{b_j(s)}
      =-\frac{\overline{\alpha_j}}r+
          \frac{r^2-|\alpha_j|^2}{r(s-\rho_j)}
      =A_j+\frac{D_j}{s-\rho_j},
                    \qquad |A_j|\leq1,\quad |D_j|\leq3.
                                                               \tag{FB.14}
\]
On the critical line the prescribed inverse transform
of its pole term is a constant of modulus at most three
times $1_{t<0}e^{-i\gamma_jt}$.
Convolving these $N$ elementary kernels gives
$A\delta_0+R_T$, with $|A|\leq1$, $\supp R_T\subset(-\infty,0]$,
and, for $u>0$,
\[
 |R_T(-u)|\leq
       \sum_{j=1}^N {N\choose j}3^j
                              \frac{u^{j-1}}{(j-1)!}.         \tag{FB.15}
\]
To see the factorial, the convolution of $j$ half-line
exponentials is an integral over the simplex of
$j$ nonnegative lengths with sum $u$, whose volume is
$u^{j-1}/(j-1)!$.  Its phases all have modulus one.
Expansion of the product gives the binomial coefficient.
This proves the bound even for coincident or arbitrarily
close ordinates.
The distribution in (FB.2) is exactly the convolution
\[
                          \kappa_T=A k_T+R_T*k_T.             \tag{FB.16}
\]

\proofstep{4. The positive-tail moments are polynomial in height.}
Fix $A_0\geq0$ and $\ell\leq k+1$.
By (FB.15), Tonelli's theorem and $r=t+u$,
the weighted $L^1$ norm on $t\geq-A_0$ of a $j$-fold
term in (FB.16) is bounded by
\[
 \frac{{N\choose j}3^j}{j!}
     \int_{r\geq-A_0}(1+A_0+r_+)^{J+j}
                                   |k_T^{(\ell)}(r)|\,dr.
                                                               \tag{FB.17}
\]
Indeed for fixed $r$, integration over
$-A_0\leq t\leq r$ integrates
$(r-t)^{j-1}$ to $(r+A_0)^j/j$; the extra weight is
at most $(1+A_0+r_+)^J$.
Equation (FB.13) applies since $J+j+1\leq M$.
After absorbing fixed powers, the sum is at most
\[
 C H^c M^C
       \sum_{j=0}^N{N\choose j}
                        C_3^j\frac{(C_4M)^j}{j!}
 \leq C H^c M^C(1+C_3)^N e^{C_4M}
 \leq C'H^{c'}.                                      \tag{FB.18}
\]
Both $N$ and $M$ are $O_{\chi,J,k}(\log H)$.
This is the step for which retaining the simplex
factorial matters.
It gives the required weighted $L^1$ bounds for
$\ell\leq k+1$.
The kernels and their derivatives vanish as $t\to+\infty$;
therefore
$\|\kappa_T^{(\ell)}\|_{L^\infty([-A_0,\infty))}
\leq\|\kappa_T^{(\ell+1)}\|_{L^1([-A_0,\infty))}$.
Multiplying this bound by the weighted $L^1$ estimate
proves the weighted $L^2$ assertion (FB.3).

\proofstep{5. Cancellation of the opposite tail and the finite window.}
The function $R_T$ on the negative half-line is a finite
exponential polynomial with frequencies $\gamma_j$ and
degree strictly below their respective multiplicities.
One way to see this without a residue division is to
successively solve the first-order convolution equations
for its elementary factors in (FB.14).
Let $R_T^{\rm ext}$ be that exponential polynomial on
the whole real line.
The same ordered integrals, with their orientation
reversed when necessary, prove the bound (FB.15) with
$u$ replaced by $|t|$ for all real $t$.
Consequently
\[
 P_T(t)=\int_{\mathbb R}R_T^{\rm ext}(t-v)k_T(v)\,dv
                                                               \tag{FB.19}
\]
is a well-defined exponential polynomial with exactly
those allowed frequencies and degrees.
For $t<0$, subtraction in (FB.16) gives
\[
 \kappa_T(t)-P_T(t)
       =A k_T(t)-\int_{v<t}R_T^{\rm ext}(t-v)k_T(v)\,dv.
                                                               \tag{FB.20}
\]
Reflection of the argument in (FB.17)--(FB.18) shows that
this remainder and its derivatives have polynomially
bounded weighted $L^1$ and $L^2$ norms on the negative
half-line as well.

Convolution by $P_T$ annihilates $W_L$, by its exact
zero-jet orthogonality.  On $t<-L$, every kernel argument
$t-s$, $s\in[-L,L]$, is negative, so one may replace
$\kappa_T(t-s)$ by the remainder in (FB.20) before
right-composing with $W_L$.
Its squared Hilbert--Schmidt norm is bounded by the corresponding
quadrant integral
$\int_0^\infty u\,|\kappa_T(-u)-P_T(-u)|^2du$.
On $t>L$ the identical assertion uses the positive
tail from (FB.3).  In these two applications use weight
$\max(J,1)$ in (FB.17); the twelve spare derivatives in
the choice of $M$ allow this also when $J=0$.
Both integrals are polynomial in $H$.

On the middle square $|t|,|s|\leq L$, (FB.13)--(FB.16)
give directly, for each fixed $\ell\leq k$,
\[
 |\kappa_T^{(\ell)}(t-s)|
                     \leq C H^c(1+2L)^N.
                                                               \tag{FB.21}
\]
For this estimate bound $|t-s-v|^{j-1}$ by
$(1+2L)^{j-1}(1+|v|)^{j-1}$ and use (FB.13);
the same factorial sum as in (FB.18) absorbs all constants.
The area of the square is $4L^2$, and right composition
by the contraction $W_L$ does not increase a
Hilbert--Schmidt norm.
These three regions prove (FB.4), after absorbing
$2^N$ in a power of $H$.
For the Sobolev norm apply the same bounds to the first
$\ell$ derivatives, or use the equivalent integer
Sobolev norm.  This completes the proof.
\end{proof}

\paragraph{Use in the actual cutoff and scope.}
For a fixed smooth endpoint multiplier the quadrant kernel
$(1-p(e^t))\kappa_T(t-s)1_{s\leq0}$ has polynomial
Hilbert--Schmidt norms with any specified finite number
of output log derivatives and log moments, by (FB.3).
The identity
$\int_{t\geq-A_0,s\leq0}F(t-s)\,dt\,ds
=\int_{-A_0}^\infty(u+A_0)F(u)\,du$
explains the extra moment used here.
Ordinary raw derivatives insert factors $e^{-t}$;
the output is bounded away from raw zero, so these
are bounded on that support.
The analytic numerator version also covers the fixed
parity Fourier multiplier in AC-T39, on these high
frequency windows; its local complex bounds follow
from its displayed Gamma formula and Stirling's formula.
No complete prolate trace limit is asserted by this
boundary estimate alone.  Its use still requires a
quantitative approximation of the target half-line
boundary operators by actual finite theta images.
This is an analytic estimate in the Hilbert realization,
not a new finite-etale cover, motivic weight, or
Grothendieck archive assertion.
